\documentclass[11pt]{article}

\usepackage[margin=2.5cm]{geometry}
\usepackage{amsmath,amssymb,amsthm,bm}

\title{Exercise Sheet -- Conditioning \& Martingales\\[0.4em]
{\large Applications of Lecture: Conditional Expectation and (Discrete-Time) Martingales}}
\author{}
\date{}

\begin{document}
\maketitle

This sheet is organized in two parts:
\begin{itemize}
  \item \textbf{Part A -- Easy/Medium applied exercises}
  \item \textbf{Part B -- Harder applied exercises (still doable)}
\end{itemize}

The problems are designed to practice:
\begin{itemize}
  \item conditioning and conditional expectation (discrete/continuous, partitions, projection intuition),
  \item filtrations generated by observations,
  \item key properties (linearity, tower, ``taking out what is known'', Jensen),
  \item martingales/submartingales/supermartingales with the \emph{natural filtration},
  \item simple finance-flavored applications (binomial pricing, fair game, risk-neutral drift).
\end{itemize}

\bigskip
\hrule
\bigskip

% =====================================================
\section*{Part A -- Easy/Medium Applied Exercises}
% =====================================================

\begin{enumerate}

% -------------------------------------------------
\item \textbf{Conditional expectation on a finite partition (quick computation)}

Let $\Omega=\{\omega_1,\omega_2,\omega_3,\omega_4,\omega_5,\omega_6\}$ with uniform probability.
Consider the partition
\[
B_1=\{\omega_1,\omega_2\},\quad B_2=\{\omega_3,\omega_4,\omega_5\},\quad B_3=\{\omega_6\}.
\]
Let $X$ be defined by
\[
X(\omega_1)=2,\ X(\omega_2)=0,\ X(\omega_3)=1,\ X(\omega_4)=4,\ X(\omega_5)=1,\ X(\omega_6)=9.
\]
\begin{enumerate}
\item Compute $\mathbb{E}[X]$.
\item Compute $\mathbb{E}[X\mid \sigma(B_1,B_2,B_3)]$ explicitly.
\item Check the averaging property $\mathbb{E}[\mathbb{E}[X\mid\sigma(B_1,B_2,B_3)]]=\mathbb{E}[X]$.
\end{enumerate}
\emph{Hint:} $\mathbb{E}[X\mid\sigma(B_1,B_2,B_3)]$ is constant on each $B_i$.

\bigskip

% -------------------------------------------------
\item \textbf{Conditional expectation of an indicator (Bayes-style)}

Let $A,B$ be events with $\mathbb{P}(B)\in(0,1)$.
\begin{enumerate}
\item Show that $\mathbb{E}[\mathbf{1}_A\mid B]=\mathbb{P}(A\mid B)$.
\item Show that
\[
\mathbb{E}[\mathbf{1}_A\mid \sigma(B)]
=
\mathbb{P}(A\mid B)\mathbf{1}_B+\mathbb{P}(A\mid B^c)\mathbf{1}_{B^c}.
\]
\item (Applied) A credit event $D$ (default) and a rating event $R$ satisfy
$\mathbb{P}(D)=0.02$, $\mathbb{P}(R)=0.1$, $\mathbb{P}(D\cap R)=0.015$.
Compute $\mathbb{E}[\mathbf{1}_D\mid \sigma(R)]$.
\end{enumerate}

\bigskip

% -------------------------------------------------
\item \textbf{Conditioning in a discrete model (signal and state)}

Let $Y\in\{G,B\}$ (Good/Bad) with $\mathbb{P}(Y=G)=0.7$ and $\mathbb{P}(Y=B)=0.3$.
A signal $S\in\{0,1\}$ satisfies
\[
\mathbb{P}(S=1\mid Y=G)=0.9,\qquad \mathbb{P}(S=1\mid Y=B)=0.2.
\]
Define the payoff $X=\mathbf{1}_{\{Y=G\}}$.
\begin{enumerate}
\item Compute $\mathbb{P}(S=1)$.
\item Compute $\mathbb{E}[X\mid S=1]$ and $\mathbb{E}[X\mid S=0]$.
\item Write $\mathbb{E}[X\mid \sigma(S)]$ as a function of $S$.
\end{enumerate}
\emph{Hint:} $\mathbb{E}[X\mid S=s]=\mathbb{P}(Y=G\mid S=s)$.

\bigskip

% -------------------------------------------------
\item \textbf{Continuous conditional density (simple and explicit)}

Let $(X,Y)$ have joint density on $\{0\le y\le x\le 1\}$:
\[
f_{X,Y}(x,y)=2\,\mathbf{1}_{\{0\le y\le x\le 1\}}.
\]
\begin{enumerate}
\item Compute $f_X(x)$.
\item Find the conditional density $f_{Y\mid X=x}(y)$.
\item Compute $\mathbb{E}[Y\mid X=x]$ and then $\mathbb{E}[Y\mid X]$.
\end{enumerate}

\bigskip

% -------------------------------------------------
\item \textbf{Tower property with a natural filtration}

Let $(X_1,X_2,X_3)$ be integrable random variables and define the natural filtration
\[
\mathcal{F}_t=\sigma(X_1,\dots,X_t),\qquad t=1,2,3.
\]
Let $Z=X_1+2X_2+3X_3$.
\begin{enumerate}
\item Compute $\mathbb{E}[Z\mid \mathcal{F}_2]$ in terms of $X_1,X_2$ and $\mathbb{E}[X_3\mid \mathcal{F}_2]$.
\item Show that $\mathbb{E}[\mathbb{E}[Z\mid \mathcal{F}_2]\mid \mathcal{F}_1]=\mathbb{E}[Z\mid \mathcal{F}_1]$.
\item (Applied) Interpret $\mathbb{E}[Z\mid \mathcal{F}_2]$ as a forecast of $Z$ given data up to time $2$.
\end{enumerate}

\bigskip

% -------------------------------------------------
\item \textbf{``Taking out what is known'' (scaling by information)}

Let $X\in L^1$ and $Y$ be a random variable. Let $g:\mathbb{R}\to\mathbb{R}$ with $\mathbb{E}[|g(Y)X|]<\infty$.
\begin{enumerate}
\item Prove that $\mathbb{E}[g(Y)X\mid Y]=g(Y)\,\mathbb{E}[X\mid Y]$.
\item (Applied) Let $R$ be a random return and let $Y$ be a ``volatility regime'' indicator. Explain why
$\mathbb{E}[\text{position}(Y)\cdot R\mid Y]=\text{position}(Y)\,\mathbb{E}[R\mid Y]$.
\end{enumerate}

\bigskip

% -------------------------------------------------
\item \textbf{Conditional variance decomposition (regression intuition)}

Let $X\in L^2$ and $Y$ a random variable.
\begin{enumerate}
\item Prove
\[
\mathrm{Var}(X)=\mathbb{E}[\mathrm{Var}(X\mid Y)]+\mathrm{Var}(\mathbb{E}[X\mid Y]).
\]
\end{enumerate}

\bigskip

% -------------------------------------------------
\item \textbf{A simple martingale from i.i.d.\ increments (fair game)}

Let $(\varepsilon_t)_{t\ge1}$ be i.i.d.\ with $\mathbb{E}[\varepsilon_t]=0$ and define
\[
S_t=\sum_{k=1}^t \varepsilon_k,\qquad \mathcal{F}_t=\sigma(\varepsilon_1,\dots,\varepsilon_t).
\]
\begin{enumerate}
\item Show that $(S_t)_{t\ge0}$ (with $S_0=0$) is a martingale w.r.t.\ $(\mathcal{F}_t)$.
\item Show that $S_t^2-t\,\mathrm{Var}(\varepsilon_1)$ is a martingale.
\item (Finance hint) Interpret $S_t$ as cumulative \emph{unexpected} P\&L and explain ``no predictable drift''.
\end{enumerate}

\bigskip

% -------------------------------------------------
\item \textbf{Submartingale via convexity (Jensen in action)}

Let $(M_t)$ be a martingale (integrable) w.r.t.\ its natural filtration $(\mathcal{F}_t)$.
Let $g:\mathbb{R}\to\mathbb{R}$ be convex and assume $g(M_t)$ is integrable.
\begin{enumerate}
\item Show that $(g(M_t))$ is a submartingale.
\item Apply to $g(x)=x^2$ and conclude that $(M_t^2)$ is a submartingale.
\item (Applied) Explain why convex payoffs (e.g.\ options) tend to have ``positive drift'' under fair-game dynamics.
\end{enumerate}
\emph{Hint:} use conditional Jensen:
$g(\mathbb{E}[M_{t+1}\mid\mathcal{F}_t])\le \mathbb{E}[g(M_{t+1})\mid\mathcal{F}_t]$.

\bigskip

% -------------------------------------------------
% -------------------------------------------------
\item \textbf{One-step binomial pricing (mini finance) -- detailed version}

A stock price satisfies \(S_0=100\). After one period,
\[
S_1=
\begin{cases}
120 & \text{(up state)},\\
90  & \text{(down state)}.
\end{cases}
\]
The \emph{physical} (real-world) probability of the up state is \(p\in(0,1)\),
but we will price by \emph{no-arbitrage}. The risk-free rate is \(r=0\),
so the bank account satisfies \(B_0=1\) and \(B_1=1\) (no discounting).

We consider a European call option with strike \(K=100\):
\[
H=(S_1-K)_+.
\]

\begin{enumerate}
% -----------------------------
\item \textbf{Payoff table.}
Compute the option payoff \(H\) in each state and fill the table:
\[
\begin{array}{c|c|c}
\text{State} & S_1 & H=(S_1-100)_+\\ \hline
\text{Up} & 120 & \ \ ? \ \\
\text{Down} & 90 & \ \ ? \ \\
\end{array}
\]

% -----------------------------
\item \textbf{No-arbitrage and risk-neutral probability.}
A probability \(q\in(0,1)\) is called \emph{risk-neutral} if the discounted stock price
is a martingale. Since \(r=0\), this means:
\[
\mathbb{E}^q[S_1]=S_0.
\]
\begin{enumerate}
\item Write \(\mathbb{E}^q[S_1]\) explicitly in terms of \(q\).
\item Solve for \(q\) and check that \(0<q<1\).
\item Explain in one sentence why \(q\) is \emph{not} necessarily equal to the real probability \(p\).
\end{enumerate}

% -----------------------------
\item \textbf{No-arbitrage option price.}
The no-arbitrage price is the risk-neutral expectation (here no discounting):
\[
V_0=\mathbb{E}^q[H].
\]
\begin{enumerate}
\item Write \(\mathbb{E}^q[H]\) using the two-state payoff table.
\item Compute \(V_0\) numerically.
\end{enumerate}

% -----------------------------
\item \textbf{Replicating portfolio.}
We look for a self-financing portfolio made of:
\begin{itemize}
  \item \(\Delta\) shares of stock,
  \item \(B\) units in the bank account (cash),
\end{itemize}
such that its terminal value matches the option payoff in both states:
\[
\Delta S_1 + B = H.
\]
\begin{enumerate}
\item Write the two equations obtained by matching the payoff in the up and down states.
\item Solve explicitly for \(\Delta\) and \(B\).
\item Compute the initial cost of the portfolio:
\[
\Delta S_0 + B,
\]
and check that it matches the price \(V_0\) found in (c).
\end{enumerate}

% -----------------------------
\item \textbf{Interpretation (short).}
\begin{enumerate}
\item Is \(B\) positive or negative? Interpret economically (lending vs borrowing).
\item Is \(\Delta\) between 0 and 1? Interpret as the “hedge ratio”.
\end{enumerate}

\end{enumerate}

\emph{Hint:} from payoff matching you should obtain
\[
\Delta\cdot 120 + B = 20,
\qquad
\Delta\cdot 90 + B = 0.
\]
\end{enumerate}

\bigskip
\hrule
\bigskip

% =====================================================
\section*{Part B -- Harder Applied Exercises}
% =====================================================

\begin{enumerate}
\setcounter{enumi}{10}

% -------------------------------------------------
\item \textbf{Conditional Jensen inequality (proof + applied bound)}

Let $X\in L^1$, $\mathcal{G}\subseteq\mathcal{F}$ a sub-$\sigma$-algebra, and $\varphi$ convex.
\begin{enumerate}
\item Prove the conditional Jensen inequality:
\[
\varphi(\mathbb{E}[X\mid \mathcal{G}])\le \mathbb{E}[\varphi(X)\mid \mathcal{G}]
\quad\text{a.s.}
\]
\item Apply it with $\varphi(x)=x^2$ and interpret the result as ``conditioning reduces MSE''.
\item (Finance) Let $R$ be a return and $\mathcal{G}$ represent information at time $t$.
Explain why $\mathbb{E}[R\mid\mathcal{G}]$ is the best predictor in mean-square error.
\end{enumerate}

\bigskip

% -------------------------------------------------
\item \textbf{Projection property (compute best linear predictor)}

Let $X,Y\in L^2$. Consider predictors $\hat X=aY+b$.
\begin{enumerate}
\item Show that the minimizer of $\mathbb{E}[(X-aY-b)^2]$ satisfies
\[
a=\frac{\mathrm{Cov}(X,Y)}{\mathrm{Var}(Y)},\qquad b=\mathbb{E}[X]-a\,\mathbb{E}[Y],
\]
assuming $\mathrm{Var}(Y)>0$.
\item Compute the minimal MSE in terms of $\mathrm{Var}(X)$ and $\rho=\mathrm{Corr}(X,Y)$.
\item (Applied) Interpret what happens when $|\rho|\approx 1$ vs.\ $\rho\approx 0$.
\end{enumerate}
\emph{Hint:} expand the square and differentiate in $a,b$.

\bigskip

% -------------------------------------------------
\item \textbf{Gaussian conditioning (explicit formulas)}

Let $(X,Y)$ be jointly Gaussian with
\[
\mathbb{E}[X]=\mu_X,\quad \mathbb{E}[Y]=\mu_Y,\quad
\mathrm{Var}(X)=\sigma_X^2,\quad \mathrm{Var}(Y)=\sigma_Y^2,\quad
\mathrm{Cov}(X,Y)=\rho\sigma_X\sigma_Y.
\]
\begin{enumerate}
\item Show that
\[
\mathbb{E}[X\mid Y]=\mu_X+\frac{\mathrm{Cov}(X,Y)}{\mathrm{Var}(Y)}(Y-\mu_Y).
\]
\item Show that
\[
\mathrm{Var}(X\mid Y)=\sigma_X^2(1-\rho^2).
\]
\item (Finance) Interpret $\rho$ as ``hedge effectiveness'' and relate $\mathrm{Var}(X\mid Y)$ to residual risk.
\end{enumerate}

\bigskip

% -------------------------------------------------
\item \textbf{Two-step binomial model (pricing by martingale)}

Let $S_0>0$, and in each period the stock moves by factors $u>1$ or $d\in(0,1)$.
Let the risk-free rate be $r> -1$.
\begin{enumerate}
\item Derive the risk-neutral probability $q$ such that the discounted price
$\tilde S_t=S_t/(1+r)^t$ is a martingale.
\item Write the time-0 price of a payoff $H=h(S_2)$ as $V_0=\mathbb{E}^q[\tilde H]$ with $\tilde H=H/(1+r)^2$.
\item Compute $V_0$ explicitly for a digital payoff $H=\mathbf{1}_{\{S_2\ge K\}}$.
\end{enumerate}
\emph{Hint:} enumerate the 3 possible values of $S_2$ and their $q$-probabilities.

\bigskip

% -------------------------------------------------
\item \textbf{Stopping time with natural filtration (optional stopping, simple case)}

Let $(S_t)$ be the simple symmetric random walk: $S_0=0$, $S_t=\sum_{k=1}^t \varepsilon_k$
with $\varepsilon_k\in\{-1,+1\}$ i.i.d.\ and $\mathbb{P}(\varepsilon_k=1)=1/2$.
Let $\mathcal{F}_t=\sigma(\varepsilon_1,\dots,\varepsilon_t)$.
Define the hitting time
\[
\tau=\inf\{t\ge 0:\ S_t\in\{-a,+b\}\},
\quad a,b\in\mathbb{N}.
\]
\begin{enumerate}
\item Argue that $\tau$ is a stopping time w.r.t.\ $(\mathcal{F}_t)$.
\item Using that $(S_t)$ is a martingale, derive
\[
\mathbb{E}[S_\tau]=0
\quad\Rightarrow\quad
\mathbb{P}(S_\tau=+b)=\frac{a}{a+b}.
\]
\item (Applied) Interpret $\frac{a}{a+b}$ as a ``ruin probability'' in a fair game.
\end{enumerate}
\emph{Hint:} $S_\tau$ only takes values $-a$ and $+b$.

\bigskip

% -------------------------------------------------
\item \textbf{Doob decomposition in a toy model (applied)}

Let $(X_t)_{t\ge0}$ be integrable and adapted to $\mathcal{F}_t=\sigma(X_0,\dots,X_t)$.
Define
\[
A_{t+1}-A_t=\mathbb{E}[X_{t+1}-X_t\mid \mathcal{F}_t],
\qquad A_0=0,
\]
and $M_t=X_t-A_t$.
\begin{enumerate}
\item Show that $(M_t)$ is a martingale.
\item Interpret $A_t$ as ``predictable drift'' and $M_t$ as ``surprise''.
\item (Finance) In a P\&L process, explain how this separates expected drift from innovation.
\end{enumerate}

\bigskip

% -------------------------------------------------
\item \textbf{Martingale vs.\ submartingale from interest accrual (finance-flavored)}

Let $V_t$ be a portfolio value process adapted to $\mathcal{F}_t=\sigma(\text{market info up to }t)$.
Assume
\[
\mathbb{E}[V_{t+1}\mid \mathcal{F}_t]\ge (1+r)V_t
\]
for a constant $r>-1$.
\begin{enumerate}
\item Show that the discounted process $\tilde V_t=V_t/(1+r)^t$ is a submartingale.
\item Give a financial interpretation (``positive expected alphaa'' after discounting).
\item What changes if the inequality is reversed?
\end{enumerate}

\bigskip

% -------------------------------------------------
\item \textbf{Pricing with conditioning (tower property in practice)}

Let $(\mathcal{F}_t)$ be the natural filtration of the market.
A discounted payoff $\tilde H$ is integrable and the (risk-neutral) price process is
\[
\tilde V_t=\mathbb{E}[\tilde H\mid \mathcal{F}_t].
\]
\begin{enumerate}
\item Show that $(\tilde V_t)$ is a martingale.
\item Show the \emph{dynamic pricing identity}:
$\tilde V_s=\mathbb{E}[\tilde V_t\mid \mathcal{F}_s]$ for $s\le t$.
\item (Applied) Explain why this justifies ``rolling'' forecasts/prices as information arrives.
\end{enumerate}

\bigskip

% -------------------------------------------------
\item \textbf{A simple risk measure and conditional expectation (applied)}

Let $L$ be a loss (positive = loss) with $L\in L^1$ and let $\mathcal{G}$ be information.
Define the conditional expected loss $\rho_\mathcal{G}(L)=\mathbb{E}[L\mid \mathcal{G}]$.
\begin{enumerate}
\item Show that if $L\le L'$ a.s., then $\rho_\mathcal{G}(L)\le \rho_\mathcal{G}(L')$ a.s.
\item Show translation invariance w.r.t.\ $\mathcal{G}$-measurable cash $C$:
$\rho_\mathcal{G}(L+C)=\rho_\mathcal{G}(L)+C$.
\item (Finance) Interpret $\mathbb{E}[L\mid \mathcal{G}]$ as a state-dependent capital requirement.
\end{enumerate}

\bigskip

% -------------------------------------------------
\item \textbf{Applied ``martingale transform'' (hedging intuition)}

Let $(M_t)$ be a martingale w.r.t.\ $\mathcal{F}_t$ and let $(H_t)$ be a bounded predictable process,
meaning $H_t$ is $\mathcal{F}_{t-1}$-measurable.
Define
\[
G_t=\sum_{k=1}^t H_k(M_k-M_{k-1}),\qquad G_0=0.
\]
\begin{enumerate}
\item Show that $(G_t)$ is a martingale.
\item (Finance) Interpret $H_t$ as a trading strategy and $G_t$ as cumulative gains.
\item In one sentence: why is predictability ($H_t$ known at time $t-1$) crucial?
\end{enumerate}

\end{enumerate}

\end{document}